\subsection{Student SID: 540840166}

\noindent 

In this task, we mainly used and compared three models: RF, MLP, and CNN. I have deepened my understanding of the principles and applicable scenarios of different models, and mastered the complete machine learning modeling process, including data preprocessing, model training, hyperparameter adjustment, and model evaluation. By comparing the performance of the three methods, I have further improved my ability to analyze problems and select appropriate models. Random forest is suitable for structured data, with simple modeling and stable results. However, it performs generally well for multidimensional image data in this task. MLP can handle nonlinear problems and is suitable for small and medium-sized numerical data. MLP performs well in this task. CNN performs particularly well in data with spatial structures such as images and is a suitable model for this task.

\subsection{Student SID: 541002619}

\noindent 

My deepest realization during the process of completing this assignment is that the choice of model has a much greater impact on the task results than I had imagined before. At first, I thought MLP or Random Forest might also be able to handle image classification problems, but after actual testing, I found that only CNN truly understands the "spatial structure" of images. It can automatically extract local features and construct more complex representations, with far superior accuracy and classification balance. Through these comparisons, I truly understand what it means for a model to adapt to the characteristics of the data.

